\section{Equivariant cohomology rings of \texorpdfstring{$\SX_\Psi$}{X\_Psi}}\label{sec_xpsi}
\subsection{Conventions for equivariant cohomology}\label{sec_pre_equ_coho}

Unless specified otherwise, all cohomology groups are taken with integral coefficients
in the sense of Borel equivariant cohomology.
Let $G = \mathrm{GL}_n(\C)$, let $T \subset G$ be the maximal torus of diagonal matrices,
and let $S = \C^\times$ act on the one-dimensional representation $\C_1 \cong \C$ by
the standard character of weight $1$ (loop rotation).
For a linear algebraic group $K$, we write $\Lambda_K = H_K^*(\pt; \Z)$.
Our primary coefficient rings are
\begin{equation}\label{eq:coefficient_rings}
 \Lambda_{T\times S} = \Z[t_1,\dots,t_n,\epsilon], \qquad
 \Lambda_{G\times S} = \Z[c_1,\dots,c_n,\epsilon],
\end{equation}
where $t_a\in H^2_T(\pt;\Z)$ is the equivariant first Chern class of the $a$-th coordinate line of $\C^n$ (corresponding to the character $\epsilon_a$), $\epsilon = c_1^S(\C_1)$, and $c_j = c_j^G(\C^n)$,
with gradings $\deg t_a = \deg\epsilon = 2$ and $\deg c_j = 2j$.
The restriction homomorphism $\Lambda_{G\times S}\hookrightarrow\Lambda_{T\times S}$ identifies
$c_j$ with the $j$-th elementary symmetric polynomial $e_j(t_1,\dots,t_n)$.
For an equivariant vector bundle $E$, we denote its equivariant Chern polynomial by
$c_u^K(E) = \sum_{j \ge 0} c_j^K(E) u^j$; this extends by multiplicativity to virtual bundles
as a formal power series in $u$.

\begin{lemma}\label{lem:borel}
Let $A$ be a complex algebraic torus (possibly trivial), and let $X$ be a
smooth complex quasi-projective variety equipped with an algebraic
$G\times A$-action.
Assume that the ordinary cohomology $H^*(X;\Z)$ is free over $\Z$ and
concentrated in even degrees.
Then the restriction homomorphism induces isomorphisms of graded algebras
\[
 H^*_{G\times A}(X;\Z)
 \otimes_{\Lambda_{G\times A}}\Lambda_{T\times A}
 \xrightarrow{\sim}H^*_{T\times A}(X;\Z),
 \qquad
 H^*_{G\times A}(X;\Z)
 \xrightarrow{\sim}H^*_{T\times A}(X;\Z)^W.
\]
Here the Weyl group $W = \mathfrak S_n$ acts on $H^*_{T\times A}(X;\Z)$ via the normalizer $N_G(T)$,
permuting the classes $t_a$, fixing $\Lambda_A$, and fixing all classes in the image of $H^*_{G\times A}(X;\Z)$.
\end{lemma}

\begin{proof}
Set $K = G \times A$ and $K_T = T \times A$.
Since $H^*(X;\Z)$ and $H^*(BK;\Z) \cong \Lambda_K$ are free $\Z$-modules
concentrated in even degrees, the Serre spectral sequence for the Borel fibration
$X \to EK \times_K X \to BK$ collapses at the $E_2$-page.
By the Leray--Hirsch theorem, any homogeneous integral basis $\{b_j\}$ of $H^*(X;\Z)$
lifts to a free $\Lambda_K$-basis $\{\widetilde b_j\}$ of $H_K^*(X;\Z)$.
These lifts restrict to a $\Lambda_{K_T}$-basis of $H_{K_T}^*(X;\Z)$,
yielding the base-change isomorphism.

Since the action of $W$ on $EK \times_{K_T} X$ covers the identity map on $EK \times_K X$,
each restricted basis element $\operatorname{res}(\widetilde b_j)$ is $W$-invariant.
Expanding an arbitrary class $\alpha \in H_{K_T}^*(X;\Z)$ in this basis, we see that $\alpha$
is $W$-invariant if and only if its coefficients lie in
$\Lambda_{K_T}^W = \operatorname{res}(\Lambda_K)$.
This identifies $H_K^*(X;\Z)$ with the invariant subalgebra $H_{K_T}^*(X;\Z)^W$.
\end{proof}

We adopt the convention that the projective bundle $\PP(E)$ parametrizes one-dimensional
linear subspaces (lines) in a vector bundle $E$ of rank $r$ over $X$.
Setting $x = -c_1^K(\CO_{\PP(E)}(-1))$, the equivariant projective bundle formula gives
\begin{equation}\label{eq:projective_bundle_formula}
 H_K^*(\PP(E);\Z) \cong
 H_K^*(X;\Z)[x]\Big/\Big(\sum_{j=0}^r c_j^K(E)x^{r-j}\Big).
\end{equation}
The classes $1, x, \dots, x^{r-1}$ form a free basis of $H_K^*(\PP(E);\Z)$ as an
$H_K^*(X;\Z)$-module, and the fiber integration map $\pi_*:H_K^*(\PP(E);\Z)\to H_K^{*-2(r-1)}(X;\Z)$
is normalized by $\pi_*(x^{r-1}) = 1$ and $\pi_*(x^i) = 0$ for $0\leq i < r-1$
(cf.\ \cite[Proposition~2.6.1]{anderson2024equivariant}).

\subsection{Lattice flags}\label{ss:def_xpsi}

Let $\CO=\C\llbracket t\rrbracket$ and $\CK=\C(\!(t)\!)$.
A \emph{lattice} in $\CK^n$ is a free $\CO$-submodule of rank $n$.
The set of all lattices in $\CK^n$ is parametrized by the \emph{affine Grassmannian}
$\Gr=G(\CK)/G(\CO)$ of $G=\mathrm{GL}_n$.

Let $S=\C^\times$. The semidirect product $G(\CO)\rtimes S$ acts on $\CK^n$ by
\[
  (g,a)\cdot v(t) = g \cdot v(at)
\]
for $g\in G(\CO)$, $a\in S$, and $v(t)\in \CK^n$, inducing an action of $G(\CO)\rtimes S$ on $\Gr$.

Following Kato \cite{kato2023geometric} and Kamnitzer \cite{kamnitzer2024fixed},
the \emph{Kato Catalan variety} $\SX_{\rfunc}$ associated with a reverse
Hessenberg function $\rfunc$ is defined by
\[
 \SX_{\rfunc}=\left\{L_\bullet=(L_0,\dots,L_n)\ \middle|\
 \begin{gathered}
 t\CO^n=L_0\subset L_1\subset\cdots\subset L_n,\quad L_i\in\Gr,\\
 \dim_\C(L_i/L_{i-1})=1,\quad tL_i\subset L_{\rfunc(i)}
 \quad(1\leq i\leq n)
 \end{gathered}\right\}.
\]
We also write $\SX_\Psi$ for $\SX_{\rfunc}$ to emphasize the corresponding root ideal $\Psi$.
The variety $\SX_{\rfunc}$ is stable under the action of $G(\CO)\rtimes S$.
The endpoint morphism is
\[
  \mathsf m_\Psi:\SX_\Psi\longrightarrow\Gr,\qquad
  (L_0,\dots,L_n)\longmapsto t^{-1}L_n.
\]
This morphism is $G(\CO)\rtimes S$-equivariant.

\begin{remark}[The initial lattice]\label{rem:lattice_normalization}
We choose $L_0=t\CO^n$ so that the quotient
$t^{-1}L_0/L_0\cong\C^n$ carries the trivial $S$-action (trivial loop rotation).
Consequently, the tautological classes $x_i=-c_1^K(L_i/L_{i-1})$ restrict to the standard
dual flag classes on the open vector subbundle $T^*_\Psi\CB$ (see Section~\ref{sec_TX_incl}).
The alternative normalization $\widetilde L_0=\CO^n$ used in
\cite[\S2.1]{kamnitzer2024fixed} and \cite[Lemma~6.1]{cautis2018quantum}
is obtained by setting $\widetilde L_i=t^{-1}L_i$; the defining lattice inclusions
remain unchanged. Under this shift, the endpoint morphism is simply given by
$\widetilde L_\bullet\mapsto\widetilde L_n=\mathsf m_\Psi(L_\bullet)$, and the tautological lines are twisted
by the $S$-representation $\C_{-1}$ of weight $-1$:
\[
 \widetilde L_i/\widetilde L_{i-1}
 \cong(L_i/L_{i-1})\otimes\C_{-1},\qquad
 \widetilde x_i=-c_1^K(\widetilde L_i/\widetilde L_{i-1})
   =x_i+\epsilon.
\]
Accordingly, in the $\CO^n$ convention, the polynomial $f_K(z)$ below is replaced by $f_K(z-\epsilon)$.
\end{remark}

We define a partial order on the set of reverse Hessenberg functions by setting
$\rfunc \leq \rfunc'$ if $\rfunc(i) \leq \rfunc'(i)$ for all $i \in [n]$.
For $\rfunc \leq \rfunc'$, the defining inclusions yield a natural
closed embedding $\SX_{\rfunc} \hookrightarrow \SX_{\rfunc'}$.

\subsection{A presentation of the equivariant cohomology ring}\label{ss:equiv_coho_xpsi}

For $K=T\times S$ or $G\times S$, set $R_K=\Lambda_K[x_1,\dots,x_n]$ and define the polynomial
\[
  f_K(z)=
  \begin{cases}
    \prod_{a=1}^n(z+t_a) & \text{for } K=T\times S,\\
    z^n+c_1z^{n-1}+\cdots+c_n & \text{for } K=G\times S.
  \end{cases}
\]
Using the operators introduced in Subsection~\ref{ss:deg_hecke_algs}, we define the ideal
\[
  I_{\rfunc,K}=(D_k(f_K)\mid1\leq k\leq n)\subset R_K.
\]
We now identify these algebraic relations with those given by the projective
bundle formula at each stage of the tower.

\begin{theorem}\label{thm:equiv_coho_xdelta}
Let $\rfunc$ be a reverse Hessenberg function and set
$f(z)=\prod_{a=1}^n(z+t_a)\in\Lambda_{T\times S}[z]$.
Then there is an isomorphism of graded $\Lambda_{T\times S}$-algebras
\[
  H^*_{T\times S}(\SX_{\rfunc};\Z) \cong \Lambda_{T\times S}[x_1,\dots,x_n]/I_{\rfunc,T\times S},
\]
where $I_{\rfunc,T\times S}=(D_k(f)\mid1\leq k\leq n)$ and the
polynomials $D_k(f)$ are defined in Subsection~\ref{ss:lemmas_polynomials}.
Under this isomorphism, $x_i$ maps to $-c_1^{T\times S}(L_i/L_{i-1})$.
\end{theorem}

\begin{proof}
Write $\Lambda=\Lambda_{T\times S}$ and suppress the equivariant superscript in the notation for Chern classes.
For each $0\leq k\leq n$, let $\SX_{\rfunc}^{[k]}$ denote the variety of partial flags $(L_0,\dots,L_k)$
satisfying the defining conditions of $\SX_{\rfunc}$ up to index $k$.
Forgetting the terminal lattice $L_k$ yields a tower of morphisms
\[
  \SX_{\rfunc}=\SX_{\rfunc}^{[n]}\longrightarrow\SX_{\rfunc}^{[n-1]}\longrightarrow\cdots
  \longrightarrow\SX_{\rfunc}^{[0]}=\pt.
\]
By \cite[Proposition~2.1]{kamnitzer2024fixed}, the $k$-th stage is
$\SX_{\rfunc}^{[k]}\cong\PP(\mathscr V_k)$, where, writing $a=\rfunc(k)$,
\[
 \mathscr V_k=t^{-1}L_a/L_{k-1},\qquad \operatorname{rk}\mathscr V_k=d_k=n-k+1+a.
\]
Its tautological line bundle is $L_k/L_{k-1}$. We also write $x_k$ for $-c_1(L_k/L_{k-1})$.
We proceed by induction on $k$ to prove that
\[
  H^*_{T\times S}(\SX_{\rfunc}^{[k]};\Z)
  \cong\Lambda[x_1,\dots,x_k]/(D_1(f),\dots,D_k(f)).
\]
The base case $k=0$ is immediate since $\SX_{\rfunc}^{[0]}=\pt$.
Assuming that the presentation holds for $\SX_{\rfunc}^{[k-1]}$, we determine the Chern
relation introduced by $\PP(\mathscr V_k)$.
From the two short exact sequences
\[
  0\longrightarrow L_{k-1}/L_0\longrightarrow t^{-1}L_a/L_0\longrightarrow\mathscr V_k\longrightarrow0,
\]
\[
  0\longrightarrow t^{-1}L_0/L_0\longrightarrow t^{-1}L_a/L_0
    \longrightarrow t^{-1}L_a/t^{-1}L_0\longrightarrow0,
\]
we obtain the following identity in equivariant $K$-theory:
\[
  [\mathscr V_k]
   =[\C^n]+\sum_{i=1}^{a}[t^{-1}L_i/t^{-1}L_{i-1}]
            -\sum_{i=1}^{k-1}[L_i/L_{i-1}].
\]
Here $\C^n=t^{-1}L_0/L_0$ is the standard representation of $T$ with trivial
$S$-action, having equivariant weights $t_1,\dots,t_n$.
Multiplication by $t^{-1}$ twists each tautological line by the $S$-representation $\C_{-1}$ of weight $-1$ (whose equivariant first Chern class is $-\epsilon$),
so $c_1(t^{-1}L_i/t^{-1}L_{i-1})=-(x_i+\epsilon)$.
The Whitney sum formula therefore gives
\begin{equation}\label{eq:chern_polynomial_Vk}
  c_u(\mathscr V_k)
  =\prod_{b=1}^n(1+t_bu)\,
    \frac{\prod_{i=1}^{a}(1-(x_i+\epsilon)u)}
         {\prod_{i=1}^{k-1}(1-x_i u)}
\end{equation}
in $H^*_{T\times S}(\SX_{\rfunc}^{[k-1]};\Z)\llbracket u\rrbracket$.

Since $F(u)=\prod_{b=1}^n(1+t_bu)$, Lemma~\ref{lem:operator_chern_relation}(ii)
with $j=a=\rfunc(k)$, together with the induction hypothesis, gives
\[
  \sum_{j=0}^{d_k}c_j(\mathscr V_k)u^j
   =\widehat P_{k,a}(u)
   =u^{d_k}D_k(f)\big|_{x_k=u^{-1}}.
\]
Comparing coefficients of $u$ yields
\[
  D_k(f)=\sum_{m=0}^{d_k}c_m(\mathscr V_k)x_k^{d_k-m}
\]
over $H^*_{T\times S}(\SX_{\rfunc}^{[k-1]};\Z)$.
By the equivariant projective bundle formula \eqref{eq:projective_bundle_formula},
the relation imposed at stage $k$ is precisely $D_k(f)=0$, completing the induction.
\end{proof}

The vector bundle sequences in the proof above also underlie the $K$-theoretic
presentation for the maximal function $\rfunc(k)=k-1$ in
\cite[Lemmas~8.13--8.14]{cautis2018quantum}.

Since the ordinary projective bundle formula shows that $H^*(\SX_{\rfunc};\Z)$
is a free $\Z$-module concentrated in even degrees, Lemma~\ref{lem:borel} applies.

\begin{corollary}\label{cor:equiv_coho_xdelta_G}
There are canonical isomorphisms
\[
  H^*_{G\times S}(\SX_{\rfunc};\Z)
  \cong H^*_{T\times S}(\SX_{\rfunc};\Z)^W
  \cong R_{G\times S}/I_{\rfunc,G\times S},
\]
with $x_i=-c_1^{G\times S}(L_i/L_{i-1})$.
Here the normalizer action of the Weyl group $W$ permutes the $t_i$ and fixes all
$x_i$ and $\epsilon$.
\end{corollary}

\begin{proof}
The first isomorphism follows from Lemma~\ref{lem:borel}.
For the second, the projective bundle tower in the proof of
Theorem~\ref{thm:equiv_coho_xdelta} is $G\times S$-equivariant.
Writing $a=\rfunc(k)$, the same exact sequences give
\[
 c_u^{G\times S}(\mathscr V_k)
 =(1+c_1u+\cdots+c_nu^n)
 \frac{\prod_{i=1}^{a}(1-(x_i+\epsilon)u)}
      {\prod_{i=1}^{k-1}(1-x_i u)}.
\]
Assuming the presentation through stage $k-1$,
Lemma~\ref{lem:operator_chern_relation}(ii) gives
\[
 c_u^{G\times S}(\mathscr V_k)
 =u^{d_k}D_k(f_{G\times S})\big|_{x_k=u^{-1}},
 \qquad d_k=n-k+1+a,
\]
since $u^nf_{G\times S}(u^{-1})=1+c_1u+\cdots+c_nu^n$.
Comparing coefficients identifies $D_k(f_{G\times S})$ with the projective
bundle relation. Starting from $\SX_{\rfunc}^{[0]}=\pt$, we therefore obtain
the presentation by induction on $k$ using the equivariant projective bundle formula.
The stated Weyl group action follows from Lemma~\ref{lem:borel} and the
$G\times S$-equivariance of the tautological lines.
\end{proof}

\begin{corollary}\label{cor:coho_xdelta_R}
With integer coefficients, restriction to a subtorus $T'\subset T\times S$ (possibly trivial) is given by base change:
\[
 H^*_{T'}(\SX_{\rfunc};\Z)
 \cong H^*_{T\times S}(\SX_{\rfunc};\Z)
       \otimes_{\Lambda_{T\times S}}\Lambda_{T'}.
\]
In particular, setting $f(z)=\prod_{a=1}^n(z+t_a)$, one has
\[
 H_T^*(\SX_{\rfunc};\Z)
 \cong\frac{\Z[t_1,\dots,t_n,x_1,\dots,x_n]}
 {\bigl(\partial_{k-1}\cdots\partial_{\rfunc(k)+1}
                      f(x_{\rfunc(k)+1})\mid1\leq k\leq n\bigr)}.
\]
The ordinary cohomology ring $H^*(\SX_{\rfunc};\Z)$ is obtained from this $T$-equivariant presentation by setting all $t_a=0$.
For $K=T\times S$ or $G\times S$ and every commutative ring $A$, one also has
\[
 H_K^*(\SX_{\rfunc};A)
 \cong H_K^*(\SX_{\rfunc};\Z)\otimes_\Z A.
\]
\end{corollary}

\begin{proof}
For base change, choose homogeneous equivariant lifts of an integral basis of ordinary cohomology.
Under restriction to $T'$, these classes continue to restrict to the same ordinary basis;
applying the Leray--Hirsch theorem to both Borel fibrations identifies the restriction with the stated tensor product.
Setting $\epsilon=0$ replaces each operator $\BS_i$ by the transposition $\tau_i$,
and $\tau_j\cdots\tau_1f(x_1)=f(x_{j+1})$, giving the displayed relations.

The coefficient-change isomorphism follows from the universal coefficient
theorem, since the integral equivariant cohomology is free in each degree.
\end{proof}

\subsection{Monomial bases and equivariant integration}

For $\Psi=\Psi_{\rfunc}$, set
\[
 d_k=n-k+1+\rfunc(k),\qquad
 d=\sum_{k=1}^n(d_k-1)=\binom n2+|\Psi|.
\]

\begin{corollary}\label{cor:normal_forms}
For $K=T\times S$ or $G\times S$, the monomials
\[
 x_1^{e_1}\cdots x_n^{e_n},\qquad 0\leq e_k<d_k,
\]
form a $\Lambda_K$-basis of $H^*_K(\SX_{\rfunc};\Z)$.
Equivariant integration extracts the coefficient of
$x_1^{d_1-1}\cdots x_n^{d_n-1}$ in the expansion in this basis.
\end{corollary}

\begin{proof}
At the $k$-th stage of the projective bundle tower, the classes
$1,x_k,\dots,x_k^{d_k-1}$ form a basis over the cohomology of the preceding
stage. Iterating gives the stated $\Lambda_K$-basis.
Since integration along the $k$-th fiber sends $x_k^{d_k-1}$ to $1$,
\[
 \int_{\SX_{\rfunc}}x_1^{d_1-1}\cdots x_n^{d_n-1}=1.
\]
Every other basis monomial has degree less than $2d$ and hence has zero
integral. The formula follows by $\Lambda_K$-linearity.
\end{proof}

\subsection{Duality and correspondence classes}\label{ss:duality}

Let $X=\SX_{\rfunc}$, $d=\dim_\C X$, and $M_K=H_K^*(X;\Z)$ for $K=T\times S$ or $G\times S$.
The equivariant integration pairing
\[
  M_K\otimes_{\Lambda_K}M_K\longrightarrow\Lambda_K,
  \qquad (\alpha,\beta)\longmapsto\int_X\alpha\cup\beta,
\]
is perfect over $\Lambda_K$.
Indeed, choose homogeneous equivariant lifts $\widetilde a_i,\widetilde b_i$
of Poincar\'e-dual bases of $H^*(X;\Z)$, ordered by nondecreasing $\deg a_i$.
By the Leray--Hirsch theorem, these lifts are $\Lambda_K$-bases of $M_K$. Since
\[
 \deg\int_X\widetilde a_i\widetilde b_j=\deg a_i-\deg a_j,
\]
the pairing matrix is lower triangular with diagonal entries $1$:
negative-degree entries vanish, and degree-zero entries are the ordinary intersection numbers.
External products of these lifts likewise lift an integral basis of $H^*(X\times X;\Z)$,
so the Leray--Hirsch theorem gives the canonical K\"unneth isomorphism
\[
  M_K\otimes_{\Lambda_K}M_K
  \xrightarrow{\sim}H_K^*(X\times X;\Z).
\]

Let $\operatorname{pr}_1,\operatorname{pr}_2:X\times X\to X$ denote the projections.
A correspondence class $\gamma\in H_K^{2d+r}(X\times X;\Z)$ defines a $\Lambda_K$-linear
endomorphism of degree $r$ on $M_K$ by
\begin{equation}\label{eq:correspondence_action}
 \mathcal T_\gamma(u)=(\operatorname{pr}_1)_*
       (\gamma\cup\operatorname{pr}_2^*u).
\end{equation}
Any closed $K$-invariant subvariety $Z\subset X\times X$ of codimension $c$ determines
a correspondence class $[Z]\in H_K^{2c}(X\times X;\Z)$ via its equivariant fundamental cycle
and Poincar\'e duality on $X\times X$.
For the projection morphisms $p_{ab}:X\times X\times X\to X\times X$, the
convolution product is
\begin{equation}\label{eq:convolution_product}
 \gamma*\eta=(p_{13})_*(p_{12}^*\gamma\cup p_{23}^*\eta).
\end{equation}
By the projection formula, $\mathcal T_{\gamma*\eta}=\mathcal T_\gamma\circ\mathcal T_\eta$,
and the diagonal class $[\Delta_X]$ acts as the identity.
Set $\mathcal C_X^r=H_K^{r+2d}(X\times X;\Z)$, graded so that convolution has degree zero.
The correspondence action induces a canonical isomorphism of graded algebras
\begin{equation}\label{eq:ambient_endomorphisms}
 \mathcal C_X^\bullet\xrightarrow{\sim}
 \operatorname{End}_{\Lambda_K}^\bullet(M_K).
\end{equation}
Indeed, under the K\"unneth isomorphism it sends $\alpha\otimes\beta$ to
$u\mapsto\alpha\int_X\beta u$, which is an isomorphism by the perfectness of the pairing.

